\section{An exact kernel for exceptional rows after a fixed forest}
\label{sec:exceptional}

The forced-long-row strategy excludes instances that already admit an easy
raw cancellation. A complementary approach keeps the long rows and contracts
the short-row structure around them. The resulting theorem concerns a
\emph{fixed} choice of top-handle forest and a fixed primal forest. It does
not optimize those choices. Its advantage is an exact description of both
the remaining decision problem and the source stopping obstruction.

\subsection{The quotient and its ground convention}

Let $A=(a_{ec})$ be the nonnegative occurrence matrix of the rows retained
after the chosen top forest; here $A$ again denotes integer counts, before
any reduction modulo two. As in \cref{sec:joint}, put
$b_{ec}=\min(2,a_{ec})$. Partition its row labels into $S\sqcup L$, where
every row of $S$ is short. Usually $L$ consists of all rows with
$\sum_c b_{ec}>2$, but the theorem allows additional exceptional rows.
Write $k=|L|$ and form the graph $H[S]$ of \cref{lem:graphic}, with
column vertices $C$ and ground $\omega$.

If $H[S]$ has a cycle, its edge rows already constitute a stopping set of
$A$. Consequently full raw cancellation is impossible in the designated
positive model, irrespective of the exceptional rows. We may therefore
assume that $H[S]$ is a forest. In particular its edges have actual unit
incidences; a repeated crossing or zero short row would have been a loop.
Include isolated column vertices when forming its components. Denote the
component containing $\omega$ by $K_\omega$, and the set of all other,
\emph{free} components by $\mathcal K$.

The exceptional-row quotient has row set $L$, column set $\mathcal K$,
and entries
\begin{equation}
 q_{eK}=\min\left(2,\sum_{c\in K} b_{ec}\right),
 \qquad e\in L,\quad K\in\mathcal K.
 \label{eq:exceptional-quotient}
\end{equation}
There is no column for $K_\omega$. Saturating before aggregation gives
the same result as aggregating the original counts and then saturating.
The sums in \eqref{eq:exceptional-quotient} count occurrences, not merely
the number of incidence sites. Two unit occurrences at different vertices
of one tree contribute two.

\begin{example}[Why the grounded component is omitted]
Take
\[
 A=\begin{pmatrix}1&0\\1&2\end{pmatrix},
 \qquad S=\{\text{first row}\},\qquad
 L=\{\text{second row}\}.
\]
The first column belongs to the grounded component, so $Q=(2)$. Both
source rows form a stopping set. If the grounded component were retained
as a quotient column, the erroneous quotient $(1\;2)$ would have a unit
pivot. There is no corresponding source pivot: the ground is not an
available one-handle, and its incident short edge cannot be removed
through the occupied column.
\end{example}

\subsection{All stopping projections and their maximal lifts}

Allow the empty stopping set. For $R\subseteq L$, mark every column
vertex $c$ with $\sum_{e\in R}b_{ec}>0$, and also mark $\omega$.
In each component $K$ of $H[S]$, let $T_K(R)$ be the minimal subtree
spanning the marked vertices in $K$. Equivalently, take the union of all
paths between marked vertices. Its edge set is empty when there are fewer
than two marked vertices. Let
\[
 T(R)=\bigcup_K E\bigl(T_K(R)\bigr)\subseteq S,
\]
using the original short-row labels on these edges. Thus the construction
records source rows, not only an isomorphism type of a smaller matrix.

\begin{theorem}[Exact exceptional-row kernel and obstruction lift]
\label{thm:residual-kernel}
Assume $H[S]$ is a forest and define $Q$ by
\eqref{eq:exceptional-quotient}.
\begin{enumerate}
\item A subset $R\subseteq L$ is the exceptional-row projection of a
stopping set of $A$ if and only if $R$ is a stopping set of $Q$.
\item If $R$ is a stopping set of $Q$, then $R\cup T(R)$ is the unique
largest stopping set of $A$ whose intersection with $L$ is exactly $R$.
\item If $R_* = \operatorname{Stop}(Q)$ is the canonical largest
stopping set, then
\begin{equation}
 \operatorname{Stop}(A)=R_*\cup T(R_*).
 \label{eq:stopping-lift}
\end{equation}
In particular, $A$ is completely peelable if and only if $Q$ is
completely peelable. By the prior stopping-persistence theorem, the same
equivalence holds for existence of full positive raw cancellation.
\end{enumerate}
\end{theorem}

\begin{proof}
Let $P$ be a source stopping set and put $R=P\cap L$. Suppose a free
component $K$ has total exceptional incidence exactly one on $R$. This
incidence occurs at one vertex $v$, with coefficient one. Consider the
subforest consisting of the selected short edges $P\cap S$ in $K$. If
$v$ is isolated in it, the source column sum at $v$ is one, a
contradiction. Otherwise its selected tree has at least two leaves. At
most one is $v$, and none is ground. A different leaf has one selected
short incidence and no exceptional incidence, again a contradiction.
Thus no free component has total exceptional incidence one. Saturation
preserves this predicate, so $R$ is a stopping set of $Q$.

We also claim $P\cap S\subseteq T(R)$. Each nonground leaf of a
nontrivial selected short tree must have a positive incidence from $R$;
otherwise its source column sum would be one. Every leaf is therefore
marked, with ground allowed as an additional marked leaf. Every selected
edge lies on a path between two leaves of its selected tree. That path
is the unique path between those vertices in the original forest, so the
edge belongs to $T(R)$. This proves the claim and the maximality assertion
once the proposed lift is shown to be stopping.

Conversely, assume $R$ is stopping for $Q$. At a nonground vertex of
$T(R)$ with degree at least two, the selected short rows alone contribute
at least two. A nonground leaf is marked and therefore receives at least
one exceptional incidence as well as its one short incidence. Its total
is again at least two. A vertex with degree zero and positive exceptional
incidence can only be the single marked vertex of a free component:
otherwise a path to another marked vertex, including ground when present,
would meet it. In this case the entire exceptional incidence of that free
component is concentrated there. The stopping condition for $Q$ makes
that total at least two. All other degree-zero column vertices receive
zero. Ground itself is not a matrix column. Hence every source column
sum of $R\cup T(R)$ is zero or at least two, proving that it is stopping.

The two directions establish the first two assertions, including empty
sets. For the third, every source stopping set projects into $R_*$, and
the marked subtree construction is monotone under inclusion of $R$.
Thus every source stopping set is contained in $R_*\cup T(R_*)$.
This set is itself stopping by the second assertion, proving
\eqref{eq:stopping-lift}. A matrix peels completely exactly when its
largest stopping set is empty. Finally, the equivalence between full
peeling and full positive raw cancellation is the prior result of
\cite{incoming-stopping}; no geometric operation beyond that positive
model is introduced here.
\end{proof}

The theorem describes a whole family of obstructions. A quotient stopping
set can have several source stopping extensions, but its largest extension
has the explicit minimal-subtree formula. The quotient therefore preserves
more information than a successful-cancellation boolean or a rank test.

\subsection{Constructive schedules and source replay}

The decision equivalence has a direct constructive lift. Suppose a
quotient pivot deletes exceptional row $e$ at free component $K$. Among
the currently active exceptional rows, the total incidence in $K$ is
exactly one. There is therefore a unique source incidence $(e,c)$ in
that component, and its coefficient is one. Root the current short tree
at $c$ and remove its other leaves inward. Each removed nonground leaf
has one short edge and no active exceptional incidence, so its source
column is a legal unit leaf. After these removals, $(e,c)$ is itself a
legal source pivot.
Once this quotient column is used, its unique active exceptional incidence
is deleted and no future active exceptional incidence remains in that
free component. Pruning it therefore removes no edge needed by the final
marked subtree $T(R_*)$.

Apply this procedure to each quotient pivot. If every exceptional row is
removed, the remaining short forest peels by \cref{lem:graphic}. If the
quotient stops at $R_*$, mark its incidence sites and ground, then remove
every unmarked nonground leaf of the remaining short forest. Exactly the
minimal subtrees in \eqref{eq:stopping-lift} remain. This gives a maximal
source peeling prefix and its exact canonical residual even in an
incomplete case.

The lifted trace consists only of pairs of original row and column
labels. A verifier can check each pair by summing its input column over
the then-active source rows, without trusting the quotient construction.
Checking that the final residual has no column sum one certifies
maximality. Unlike a producer using general positive pivots, this
no-fill trace retains the original entries throughout. General positive
pivots may change the labels and entries of a surviving obstruction;
their residual row list need not equal $\operatorname{Stop}(A)$.

\subsection{A bounded number of exceptional column patterns}

The quotient can still contain many free components with identical effects
on its $k$ rows. Order $L$ and associate to a quotient column $K$ its
signature
\[
 \sigma(K)=(q_{eK})_{e\in L}\in\{0,1,2\}^{k}.
\]
Delete all-zero columns and retain one representative of each nonzero
signature. Denote the resulting matrix by $\widehat Q$. Keep every row,
including zero rows, and retain the component-to-representative map for
lifting to the source.

\begin{corollary}[Column-pattern kernel]
\label{cor:pattern-kernel}
The matrices $Q$ and $\widehat Q$ have exactly the same stopping row
subsets. Therefore $\widehat Q$ determines
$\operatorname{Stop}(A)$ through \cref{thm:residual-kernel}. It has $k$
rows and at most
\[
 \min\bigl(|\mathcal K|,3^k-1\bigr)
\]
columns, with entries in $\{0,1,2\}$.
\end{corollary}

\begin{proof}
Equal column signatures have equal sums on every subset of the row set
$L$, and hence have the same predicate of being a unit column on that
subset. Repeating this predicate adds no stopping constraint. A zero
column imposes none. These observations prove preservation of the entire
stopping family. There are $3^k-1$ possible nonzero ternary signatures.
\end{proof}

Preserving zero rows is essential. A row whose incidences all lie in the
grounded component becomes a zero quotient row and is a singleton
quotient stopping set. Its source obstruction is the row together with
the appropriate paths to ground. Removing that row would change the
answer. By contrast, the single long row $(1,1,1)$ with no short rows
gives three identical unit quotient columns and compresses to $(1)$;
its successful pivot is preserved.

\subsection{Complexity and the remaining scope}

Let $N$ count the explicit row labels, column labels, and nonzero entries,
and put $s=|\mathcal K|$. Forest testing, component labeling, incidence
aggregation, and leaf pruning admit $O(N)$ graph operations. Comparing a
binary input count with two requires no expansion in its magnitude.
For a simple deterministic implementation, forming dense signatures takes
$O(ks)$ ternary entries and lexicographically sorting them takes
$O(ks\log(s+2))$ symbol comparisons. Only signatures of components
actually present are formed: there is no loop over $3^k$ possible
patterns. Sparse alternatives are available, but unnecessary for the
polynomial preprocessing claim. Label sorting and integer input handling
contribute their ordinary bit costs. The supplied research implementation
also uses direct source scans during lifting; these are polynomial and
are not presented as an optimal implementation bound.

The bounded decision matrix and the source lift data serve different
purposes. The matrix has at most $k(3^k-1)$ ternary entries; the forest and
the maps needed to reproduce original source rows can retain linear-sized
input information. No bound on $k$ in terms of crossing number, hierarchy
depth, or coordinate bit length is proved here. Nor does the theorem
control how the quotient changes under a different top or primal forest.
In particular, \cref{prop:not-matroid} still prevents replacing the
arbitrary-row joint problem by the short-row matroid argument.

For one explicitly supplied matrix, ordinary sparse column peeling
already gives a linear-time decision. The kernel is a structural reduction
that isolates the exceptional incidence data and reconstructs exact source
obstructions; it is not by itself a faster one-shot recognition algorithm.
Its potential use in repeated geometric choices requires further theorems
bounding the exceptional interface and the cost of updating it. It also
does not supply the cyclic order or boundary-mark transport needed to turn
counts into a marked geometric certificate. The construction is new to
the present repository audit relative to the inspected incoming packages;
we make no assertion of priority over the full literature on stopping
sets or graph reduction.
